%=========== ΛΥΣΗ - 9ο ΘΕΜΑ Δ ===========
\begin{thema}{Δ}
Επειδή
\[
f''(x)>0
\]
για κάθε $x\in\mathbb{R}$, η $f'$ είναι γνησίως αύξουσα στο $\mathbb{R}$ και η $f$ είναι αυστηρά κυρτή στο $\mathbb{R}$.
\begin{erwthma}
\item Έστω $x\in(0,1)$. Εφαρμόζοντας το Θεώρημα Μέσης Τιμής για την $f$ στα διαστήματα $[0,x]$ και $[x,1]$, υπάρχουν $\xi_1\in(0,x)$ και $\xi_2\in(x,1)$ τέτοια ώστε
\[
\frac{f(x)-f(0)}{x-0}=f'(\xi_1)
\]
και
\[
\frac{f(1)-f(x)}{1-x}=f'(\xi_2).
\]
Επειδή
\[
\xi_1<\xi_2
\]
και η $f'$ είναι γνησίως αύξουσα, έχουμε
\[
f'(\xi_1)<f'(\xi_2).
\]
Με χρήση των σχέσεων $f(0)=f(1)=0$, προκύπτει
\begin{align*}
\frac{f(x)}{x}
&<-\frac{f(x)}{1-x}\\
\Leftrightarrow (1-x)f(x)&<-xf(x)\\
\Leftrightarrow f(x)&<0,
\end{align*}
αφού $x(1-x)>0$. Επομένως
\[
{f(x)<0}
\]
για κάθε $x\in(0,1)$.

\item Η $f$ είναι συνεχής στο $[0,1]$, παραγωγίσιμη στο $(0,1)$ και
\[
f(0)=f(1)=0.
\]
Επομένως, σύμφωνα με το θεώρημα \eng{Rolle}, υπάρχει $x_0\in(0,1)$ τέτοιο ώστε
\[
f'(x_0)=0.
\]
Η ρίζα αυτή είναι μοναδική, διότι η $f'$ είναι γνησίως αύξουσα στο $\mathbb{R}$.

Η εφαπτομένη της $C_f$ στο σημείο με τετμημένη $x_0$ έχει συντελεστή διεύθυνσης
\[
\lambda=f'(x_0)=0.
\]
Άρα είναι παράλληλη με τον άξονα $x'x$. Συνεπώς υπάρχει ακριβώς ένα $x_0\in(0,1)$ με τη ζητούμενη ιδιότητα.

\item Επειδή η $f'$ είναι γνησίως αύξουσα και $f'(x_0)=0$, έχουμε:
\begin{itemize}
\item αν $x<x_0$, τότε $f'(x)<f'(x_0)=0$,
\item αν $x>x_0$, τότε $f'(x)>f'(x_0)=0$.
\end{itemize}
Στον παρακάτω πίνακα βλέπουμε το πρόσημο της $f'$ και τη μονοτονία της $f$:
\begin{center}
\begin{tikzpicture}
\tkzTabInit[color,lgt=2,espcl=1.8,colorC = maincolor!50!white,colorL = maincolor!20!white,colorV = maincolor!50!white]%
{$x$ / .8,$f'(x)$ /0.8,$f(x)$/0.8}%
{$-\infty$,$x_0$,$+\infty$}%
\tkzTabLine{, -, z, +, }
\tkzTabLine{,\searrow,t,\nearrow,}
\end{tikzpicture}
\end{center}

Άρα η $f$ είναι \textbf{γνησίως φθίνουσα} στο $(-\infty,x_0]$ και \textbf{γνησίως αύξουσα} στο $[x_0,+\infty)$. Επομένως παρουσιάζει μοναδικό ολικό ελάχιστο στο $x_0$, με τιμή $f(x_0)<0$.

\item Ολοκληρώνουμε κατά παράγοντες:
\begin{align*}
\int_0^1 xf''(x)\d x
&=\left[xf'(x)\right]_0^1-\int_0^1f'(x)\d x\\
&=f'(1)-\left[f(x)\right]_0^1\\
&=f'(1)-\left(f(1)-f(0)\right)\\
&=2-(0-0)\\
&={2}.
\end{align*}
\end{erwthma}
\end{thema}
%=========== ΤΕΛΟΣ ΛΥΣΗΣ - 9ο ΘΕΜΑ Δ ===========
